\metatitle{Exchange graphs}
\metadescription{Cycle exchanges, monotone paths, ordinary basis graphs, and
direct sums of colored matroids.}

\section[exchangeGraphs]{Exchange graphs}

For an ordinary matroid, single-element exchanges connect all bases. The same
phenomenon holds for colored matroids, with simple color cycles replacing
single-element exchanges.

\begin{definition}[Exchange graph]
The \defin{exchange graph} $G(\mathcal F)$ has the words of $\mathcal F$ as
vertices. Two words are adjacent if their entire disagreement is one cyclic
exchange. Equivalently, their disagreement arrows form one directed simple
cycle on the colors.
\end{definition}

Thus an edge changes the colors on one set of positions by a cyclic
permutation. With two colors, this changes one $1$ to $0$ and one $0$ to $1$,
so the exchange graph is the ordinary basis graph.

\subsection[monotonePaths]{Monotone exchange paths}

\begin{proposition}[Monotone paths]
Let $u,v\in\mathcal F$. There is a path
\[
 u=u_0,u_1,\dotsc,u_t=v
\]
in $G(\mathcal F)$ such that, at every step, $u_{j+1}$ agrees with $v$ at
strictly more positions than $u_j$ does. In particular,
\[
 t\leq
 \left\lfloor\frac{\lvert\{e\in E:u(e)\ne v(e)\}\rvert}{2}\right\rfloor.
\]
\end{proposition}

\begin{proof}
Suppose that $u_j\ne v$, and choose a position where they differ. Apply
symmetric cycle exchange to $u_j$ and $v$ through that position, and let
$u_{j+1}$ be the resulting word obtained from $u_j$. At every position in the
chosen cycle, the new entry of $u_{j+1}$ is the entry of $v$. All other entries
are unchanged. Hence no agreement is lost, and at least two disagreements are
removed. Repeating this step reaches $v$ and gives the stated bound.
\end{proof}

In particular, the exchange graph of every colored matroid is connected. The
proof gives more: once a position agrees with the destination word, the path
never changes it again.

\subsection[colorProjections]{Color projections and basis graphs}

Fix a color $c$. Recall that the sets $B(u,c)$ are the bases of the projection
matroid $M(\mathcal F,c)$. For a basis $B$ of this matroid, write
\[
 \defin{\mathcal F(c,B)}
 \coloneqq
 \{u\in\mathcal F:B(u,c)=B\}.
\]

\begin{proposition}[Lifting basis exchanges]
Suppose that $B$ and $B\setminus\{e\}\cup\{f\}$ are bases of
$M(\mathcal F,c)$. From every word $u\in\mathcal F(c,B)$, there is an edge of
$G(\mathcal F)$ to a word in
$\mathcal F(c,B\setminus\{e\}\cup\{f\})$.
\end{proposition}

\begin{proof}
Choose a word $v$ whose color-$c$ class is
$B\setminus\{e\}\cup\{f\}$. In the disagreement graph of $u$ and $v$, the
only arrow leaving color $c$ occurs at $e$, and the only arrow entering color
$c$ occurs at $f$. A cycle through $e$ must therefore contain $f$. Symmetric
cycle exchange gives a word adjacent to $u$ whose color-$c$ class is the
required basis.
\end{proof}

\begin{corollary}[Fibers and distances]
Every path in the basis graph of $M(\mathcal F,c)$ lifts from any word over
its initial basis to a path in $G(\mathcal F)$. Each fiber
$\mathcal F(c,B)$ induces a connected subgraph, and
\[
 \operatorname{dist}_{G(\mathcal F)}(u,\mathcal F(c,B))
 =\lvert B(u,c)\setminus B\rvert.
\]
\end{corollary}

\begin{proof}
Lift the edges of a basis path one at a time. If two words lie in the same
fiber, they agree at every position colored $c$. The monotone-path argument
between them therefore uses no cycle containing $c$ and stays inside the
fiber.

Each edge of $G(\mathcal F)$ changes a color class by at most one ordinary
basis exchange, giving the lower bound on the distance. Ordinary basis
exchange gives a path to $B$ of length $\lvert B(u,c)\setminus B\rvert$, and
lifting that path gives the upper bound.
\end{proof}

\subsection[directSumGraphs]{Direct sums}

The direct sum from the \href{operations.htm}{operations chapter} has a
corresponding graph operation.

\begin{proposition}[Exchange graph of a direct sum]
For colored matroids on disjoint position sets,
\[
 G(\mathcal F_1\oplus\mathcal F_2)
 \cong
 G(\mathcal F_1)\mathbin{\square}G(\mathcal F_2),
\]
where $\square$ denotes the Cartesian product of graphs.
\end{proposition}

\begin{proof}
Consider two adjacent words in the direct sum. Their disagreement arrows form
one directed simple cycle. Within either position set, the arrows are balanced
at every color because the two restrictions have the same color quotas. A
balanced subset of the edges of a directed simple cycle is either empty or the
whole cycle. Hence all positions changed by the edge lie in one factor.

Conversely, an exchange edge in either factor remains an exchange edge after
the word in the other factor is fixed. These are exactly the edges of the
Cartesian product.
\end{proof}

\begin{example}[A square]
Let
\[
 \mathcal F_1=\{\mathtt{123},\mathtt{231}\},
 \qquad
 \mathcal F_2=\{\mathtt{12},\mathtt{21}\}.
\]
Each factor has a single exchange edge. Their direct sum has exchange graph
\[
 \begin{array}{ccc}
 \mathtt{12312}&\longleftrightarrow&\mathtt{23112}\\
 \updownarrow&&\updownarrow\\
 \mathtt{12321}&\longleftrightarrow&\mathtt{23121}
 \end{array}
\]
the Cartesian product of two edges.
\end{example}
